% !TeX root = ../main.tex
\begin{frame}{器件验证与零泊松比支架}
  \begin{SlideCanvas}
    \BodyAt{55}{84}{850}{现有研究已将零泊松比结构用于支架，并结合聚合物试验、金属制造和镍钛压握模型评价长度与承载性能。}
    \PicAt{55}{157}{419}{205}{assets/zero2026_geometry.png}
    \CaptionAt{55}{376}{419}{Yasmin等：箭头单胞、平面阵列与管状结构}
    \PointAt{55}{417}{419}{器件性能的共同设计}
    \DetailAt{55}{453}{419}{几何构型、材料行为与制造方式共同影响响应。}
    \PointAt{530}{151}{370}{Wang等，2021：PCL}
    \DetailAt{530}{191}{370}{3D打印、弯曲与径向扩张试验，观察轴向长度的稳定性。}
    \PointAt{530}{264}{370}{Gupta等，2026：金属原型}
    \DetailAt{530}{304}{370}{飞秒激光切割与抛光，径向强度提高，同时出现轴向伸长。}
    \PointAt{530}{377}{370}{Yasmin等，2026：镍钛}
    \DetailAt{530}{417}{370}{压握／释放模型及打印原型。15--25°构型在50\%压握终态报告$|\mathrm{FS}|\leq0.12\%$。}
    \Sources{\RefWang；\RefGupta。\\\RefYasmin，图4、第2.3与3.5节。各项结果按其材料及加载条件解释。}
  \end{SlideCanvas}
  \BICConclusionBand{零泊松比与径向力已有共同研究，响应轨迹定制构成\BICBandEmph{进一步的设计问题}}
  \note{原有制造实验与零泊松比器件两页合并，以最直接相关的PCL、金属制造和NiTi研究作为三个节点。Gupta摘要中的径向强度3.22对1.26N/mm，短缩率-9.44\%对+2.49\%属于同研究比较，负短缩表示伸长。Yasmin的|FS|指标针对t=0.50mm、l=4.50mm、15--25度构型的50\%压握终态，不是释放全过程最大长度误差。其原型材料状态、端部约束及实际径向验证等研究范围集中在挑战页与文献说明中讨论。本页不把已发表器件工作或拟议轨迹设计视为临床结局证据。Vellaparambil2024的制造几何反馈在挑战页引用，Hasanzadeh2025保留为设计指标参考，完整文献材料不因综述压缩而删除。}
\end{frame}
